% F(AI)^2R: writers x readers --- the (AI)^2 factor unpacked into the
% four cells of human-or-AI as writer x human-or-AI as reader.
% This figure makes the central distinction of the paper visible at a
% glance: the (AI)^2 factor is squared because both writer and reader
% can be human or AI, and the disclosure surface --- what each cell
% requires from the manuscript bundle --- differs across cells.
% Loaded by paper/sections/pattern.tex via \input, immediately after
% figures/axes which renders the four FAIR axes; this figure renders
% the (AI)^2 factor itself.
\begin{figure}[!t]
\centering
\begin{tikzpicture}[
  every node/.style={font=\sffamily\footnotesize, align=center,
                     inner sep=4pt},
  cell/.style={draw, line width=0.5pt, minimum width=58mm,
               minimum height=24mm, text width=54mm,
               align=left, fill=white, anchor=center,
               draw=dlrGrau1},
  hh/.style={fill=dlrGruen5, draw=dlrGruen1, line width=0.6pt},
  ha/.style={fill=dlrBlau5,  draw=dlrBlau1, line width=0.6pt},
  ah/.style={fill=white,     draw=dlrBlau1, line width=0.6pt, dashed},
  aa/.style={fill=dlrGrau5,  draw=dlrGrau1, line width=0.6pt, dotted},
  axislbl/.style={font=\sffamily\bfseries\footnotesize, text=dlrGrau1,
                  align=center},
  axisrole/.style={font=\sffamily\footnotesize, text=dlrGrau1,
                   align=center, inner sep=2pt},
  cellhdr/.style={font=\sffamily\bfseries\footnotesize, text=black},
  cellsub/.style={font=\sffamily\itshape\footnotesize, text=dlrGrau1},
  cellnote/.style={font=\sffamily\footnotesize, text=black},
  surface/.style={font=\sffamily\footnotesize, text=dlrAccent},
]
% --- Axis labels (writer = top row, reader = left column) ----------
\node[axislbl] at (-3.6, 1.85) {writer};
\node[axisrole, anchor=south] at (-1.55, 1.55) {human};
\node[axisrole, anchor=south] at ( 1.55, 1.55) {AI};
\node[axislbl, rotate=90] at (-5.85, -0.15) {reader};
\node[axisrole, rotate=90, anchor=south] at (-5.55, 1.05)  {human};
\node[axisrole, rotate=90, anchor=south] at (-5.55, -1.45) {AI};

% Hairline rule between the two writer columns / reader rows
\draw[draw=dlrGrau3, line width=0.3pt]
  (-4.50, 1.30) -- (-4.50, -1.55);
\draw[draw=dlrGrau3, line width=0.3pt]
  (-4.95, 1.25) -- (4.65, 1.25);

% --- Cells (2x2) ---------------------------------------------------
% Top-left: human writer, human reader (the classical case)
\node[cell, hh] (HH) at (-1.55, 0.40) {%
  {\cellhdr classical scholarship}\\[1pt]
  {\cellsub all four FAIR axes}\\[2pt]
  {\cellnote prose, citations, manuscript-only;
   transcripts and prompts are out of scope.}};
% Top-right: AI writer, human reader (the paper's primary case)
\node[cell, ha] (AH) at ( 1.55, 0.40) {%
  {\cellhdr LLM-assisted writing}\\[1pt]
  {\cellsub the F(AI)\textsuperscript{2}R primary case}\\[2pt]
  {\cellnote prose + transcripts + prompts + ladder-status +
   per-claim provenance map.}};
% Bottom-left: human writer, AI reader (the audit case)
\node[cell, ah] (HA) at (-1.55, -2.20) {%
  {\cellhdr AI-assisted reading}\\[1pt]
  {\cellsub auditor / scrutinizer / search agent}\\[2pt]
  {\cellnote machine-readable claims, IRIs, SPARQL-able graph;
   licence permits ingestion.}};
% Bottom-right: AI writer, AI reader (the systemic case)
\node[cell, aa] (AA) at ( 1.55, -2.20) {%
  {\cellhdr AI-to-AI corpus}\\[1pt]
  {\cellsub training data / ML retrieval / agent-mediated review}\\[2pt]
  {\cellnote everything above, plus base-rate disclosure on the
   training-set side; collapse risk.}};

% --- Disclosure-surface labels (right edge) ------------------------
\node[surface, anchor=west] at (4.85,  0.40)
  {\textbf{disclosure surface:}\\\textit{paper bundle widens}\\\textit{top-left $\to$ bottom-right}};

% --- Caption-ready footer band -------------------------------------
\node[font=\sffamily\itshape\footnotesize, text=dlrGrau1,
      anchor=west, align=left]
  at (-4.95, -3.85)
  {The (AI) factor squared: both writer and reader can be human or AI.
   The paper's primary case is the top-right cell; its auditability
   premise is that bottom-left and bottom-right are coming next.};
\end{tikzpicture}
\caption{The (AI)\textsuperscript{2} factor unpacked as a $2\times 2$
grid of writer (columns) and reader (rows). The top-left cell
(human writer, human reader) is the classical scholarly case and
needs only the manuscript proper; the top-right cell (AI writer,
human reader) is this paper's primary case and is what every
practice in \S\ref{sec:eight} addresses; the bottom-left cell (human
writer, AI reader) is the audit / scrutinizer / agent-mediated
search case and requires machine-readable claims and a SPARQL-able
provenance graph; the bottom-right cell (AI writer, AI reader) is
the systemic case --- training-data ingestion, agent-to-agent
review, autonomous literature search --- and additionally requires
base-rate disclosure on the training-set side. The disclosure
surface widens monotonically along the diagonal: top-left needs
only prose and citations, bottom-right needs the full
F(AI)\textsuperscript{2}R bundle.}
\label{fig:ai-squared-grid}
\end{figure}
